\documentclass[../../../include/open-logic-section]{subfiles}

\begin{document}

\olfileid[pt]{sth}{z}{story}
\olsection{A história com mais pormenor}
	
Na \olref[story][approach]{sec}, citamos a descrição de Shoenfield do
processo de formação de conjuntos. Queremos agora escrever
mais alguns princípios, para tornar esta história um pouco mais precisa. Eis-os:
\begin{enumerate}
\item[] \stageshier. Cada conjunto forma-se por algum estágio.
\item[] \stagesord. Os estágios estão ordenados: alguns vêm \emph{antes}
de outros.\footnote{Assumiremos de fato---tacitamente---que os estágios
estão \emph{bem ordenados}. O que isto significa explica-se no
\olref[ordinals][]{chap}. É um pressuposto substancial. De fato,
usando uma técnica muito engenhosa devida a \citet{Scott1974}, este
pressuposto pode ser \emph{evitado} e depois \emph{deduzido}. (Isto
também explicará por que devemos pensar que existe um estágio inicial.) Não podemos aprofundar isso aqui; para mais, veja \citet{ButtonLT1}.} 
\item[] \stagesacc. Para qualquer estágio $S$, e para quaisquer conjuntos que se
formaram \emph{antes} do estágio~$S$: forma-se, no estágio~$S$, um conjunto cujos
membros são exatamente esses conjuntos. Mais nada se forma no estágio~$S$.
\end{enumerate}
Estes são princípios informais, mas poderemos usá-los para
justificar vários dos axiomas da teoria de conjuntos de Zermelo. 

(Convém uma palavra de cautela. Embora vamos apresentar alguns
axiomas completamente standard, com nomes completamente standard, os
princípios em itálico que acabamos de apresentar não têm nomes particulares
na literatura. São apenas alcunhas que esperamos serem úteis.)

\end{document}
